\section{Task 3: Recurrent Neural Network}

\begin{questionbox}
\textbf{Q1.} With RNNs, do you still have to pad your data? If yes, how?
\end{questionbox}

Yes. RNNs can process variable-length sequences conceptually, but mini-batch training still requires
rectangular tensors. Therefore, we pad sequences inside each mini-batch. Unlike the FFNN case, padding
is dynamic: each batch is padded only to the longest sequence in that batch. We also use packed
sequences, so the recurrent layer ignores padded timesteps.

\begin{questionbox}
\textbf{Q2.} Do you have to truncate the testing sequences? Justify your answer.
\end{questionbox}

No. RNNs reuse the same recurrent weights at every timestep, so they are not constrained by a fixed
input dimensionality. Test sequences longer than the longest training sequence can still be processed.
The only requirement is that each API call is mapped to a known ID or to the \texttt{<UNK>} token.

\begin{questionbox}
\textbf{Q3.} Is the RNN padding, if any, more memory efficient with respect to the FFNN's one? Why?
\end{questionbox}

Yes. The FFNN representation used a global fixed length of 90 for every sample. RNN batching can be
more memory efficient because padding is batch-local and packed sequences avoid computation on padded
timesteps. In terms of actual tokens, train sequences contain 1,224,878 API calls compared with
1,469,250 positions in the FFNN fixed representation, i.e. 83.37\% of the fixed padded size. For the
test set, the ratio is 95.92\%, because test sequences are longer on average.

\begin{table}[h!]
\centering
\begin{tabular}{lcccc}
\hline
\textbf{Split} & \textbf{Samples} & \textbf{FFNN fixed tokens} & \textbf{Actual RNN tokens} & \textbf{Ratio} \\
\hline
Train & 16325 & 1469250 & 1224878 & 0.8337 \\
Test  & 6505  & 585450  & 561584  & 0.9592 \\
\hline
\end{tabular}
\caption{Token counts for fixed FFNN padding versus actual RNN sequence lengths.}
\label{tab:task3_padding_memory}
\end{table}

\begin{questionbox}
\textbf{Q4.} Start with a simple one-directional RNN. Is your network fast as the FFNN? If not,
where do you think that time-overhead comes from?
\end{questionbox}

No. The simple one-directional RNN is slower than the FFNN. The overhead comes from the recurrent
dependency across timesteps: the hidden state at time $t$ depends on the previous hidden state, so the
sequence dimension is harder to parallelize than the matrix operations used by a FFNN.

\begin{questionbox}
\textbf{Q5.} Train and test three variations of networks: simple one-directional RNN,
bi-directional RNN, and LSTM. Report results and observations.
\end{questionbox}

We trained three sequence models: a simple one-directional RNN, a bidirectional RNN, and a
bidirectional LSTM. Hyperparameters were selected with Optuna on a stratified validation split from
the training set, optimizing macro F1. The search space is reported in
Table~\ref{tab:task3_optuna_space}.

\begin{table}[h!]
\centering
\begin{tabular}{llll}
\hline
\textbf{Hyperparameter} & \textbf{Search space} & \textbf{Scale/type} & \textbf{Models} \\
\hline
\texttt{embedding\_dim} & \{16, 32, 64\} & Categorical & All \\
\texttt{hidden\_dim} & \{64, 128, 256\} & Categorical & All \\
\texttt{num\_layers} & \{1, 2\} & Categorical & All \\
\texttt{dropout} & $[0, 0.5]$ & Uniform & All \\
\texttt{batch\_size} & \{128, 256\} & Categorical & All \\
\texttt{lr} & $[10^{-5}, 10^{-3}]$ & Log-uniform & All \\
\texttt{weight\_decay} & $[10^{-6}, 10^{-2}]$ & Log-uniform & All \\
\hline
\end{tabular}
\caption{Optuna search space for the RNN models.}
\label{tab:task3_optuna_space}
\end{table}

\begin{table}[h!]
\centering
\begin{tabular}{lccccccc}
\hline
\textbf{Model} & \textbf{Emb.} & \textbf{Hidden} & \textbf{Layers} & \textbf{Dropout} &
\textbf{LR} & \textbf{Weight decay} & \textbf{Batch} \\
\hline
Simple RNN & 64 & 128 & 1 & 0.4077 & 0.000259 & 0.000824 & 128 \\
BiRNN      & 64 & 128 & 2 & 0.4828 & 0.000414 & 0.000017 & 256 \\
BiLSTM     & 64 & 128 & 2 & 0.4828 & 0.000414 & 0.000017 & 256 \\
\hline
\end{tabular}
\caption{Best hyperparameters found for the three RNN variants.}
\label{tab:task3_best_hyperparameters}
\end{table}

\begin{figure}[h!]
\centering
\begin{minipage}{0.32\textwidth}
\centering
\includegraphics[width=\textwidth]{figures/task-3/simple_rnn_training_curve.pdf}
\end{minipage}
\hfill
\begin{minipage}{0.32\textwidth}
\centering
\includegraphics[width=\textwidth]{figures/task-3/bidirectional_rnn_training_curve.pdf}
\end{minipage}
\hfill
\begin{minipage}{0.32\textwidth}
\centering
\includegraphics[width=\textwidth]{figures/task-3/lstm_training_curve.pdf}
\end{minipage}
\caption{Validation macro F1 curves for Simple RNN (left), BiRNN (center), and BiLSTM (right).}
\label{fig:task3_training_curves}
\end{figure}

\begin{table}[h!]
\centering
\begin{tabular}{lccccc}
\hline
\textbf{Model} & \textbf{Accuracy} & \textbf{Balanced acc.} & \textbf{Macro F1} &
\textbf{F1 Goodware} & \textbf{F1 Malware} \\
\hline
Simple RNN & 0.9600 & 0.8032 & 0.7607 & 0.5423 & 0.9791 \\
BiRNN      & 0.8999 & 0.8135 & 0.6477 & 0.3497 & 0.9458 \\
BiLSTM     & 0.9545 & 0.8597 & 0.7651 & 0.5542 & 0.9760 \\
\hline
\end{tabular}
\caption{Final test performance of the three RNN variants.}
\label{tab:task3_test_performance}
\end{table}

\begin{table}[h!]
\centering
\begin{tabular}{lcccc}
\hline
\textbf{Model/Class} & \textbf{Precision} & \textbf{Recall} & \textbf{F1-score} & \textbf{Support} \\
\hline
Simple RNN Goodware (0) & 0.4738 & 0.6337 & 0.5423 & 243 \\
Simple RNN Malware (1)  & 0.9856 & 0.9727 & 0.9791 & 6262 \\
BiRNN Goodware (0)      & 0.2309 & 0.7202 & 0.3497 & 243 \\
BiRNN Malware (1)       & 0.9882 & 0.9069 & 0.9458 & 6262 \\
BiLSTM Goodware (0)     & 0.4371 & 0.7572 & 0.5542 & 243 \\
BiLSTM Malware (1)      & 0.9903 & 0.9622 & 0.9760 & 6262 \\
\hline
\end{tabular}
\caption{Per-class classification report for the RNN variants.}
\label{tab:task3_classification_report}
\end{table}

\begin{figure}[h!]
\centering
\begin{minipage}{0.32\textwidth}
\centering
\includegraphics[width=\textwidth]{figures/task-3/simple_rnn_confusion_matrix.pdf}
\end{minipage}
\hfill
\begin{minipage}{0.32\textwidth}
\centering
\includegraphics[width=\textwidth]{figures/task-3/bidirectional_rnn_confusion_matrix.pdf}
\end{minipage}
\hfill
\begin{minipage}{0.32\textwidth}
\centering
\includegraphics[width=\textwidth]{figures/task-3/lstm_confusion_matrix.pdf}
\end{minipage}
\caption{Confusion matrices for Simple RNN (left), BiRNN (center), and BiLSTM (right).}
\label{fig:task3_confusion_matrices}
\end{figure}

The best overall model was the bidirectional LSTM, with macro F1 0.7651 and goodware F1 0.5542.
The simple RNN was close in macro F1, while the bidirectional RNN had higher goodware recall but many
more false positives, which reduced its accuracy and macro F1. Qualitatively, this means that the
BiRNN was more aggressive in predicting the minority goodware class, whereas the Simple RNN and
BiLSTM kept a better balance between detecting goodware and preserving malware accuracy.

\begin{questionbox}
\textbf{Q6.} Is the RNNs training as stable as the FFNN's one?
\end{questionbox}

No. RNN training is generally less stable than FFNN training because recurrent models propagate
information through many sequential steps, making optimization more sensitive to gradients and
learning-rate choices. To analyze this properly, we compare training loss and validation loss over the
final training epochs: lower loss is better, and a large gap where training loss keeps decreasing
while validation loss increases would indicate overfitting. We also monitor validation macro F1,
because it is the model-selection metric and is more informative than accuracy on this imbalanced
dataset.

\begin{questionbox}
\textbf{Q7.} How does your model's performance compare to the simple frequency baseline, given that
you now account for the sequence of API calls and use a significantly more complex network?
\end{questionbox}

The best RNN model, the bidirectional LSTM, slightly outperformed the frequency-based Logistic
Regression baseline in macro F1 and goodware F1. Logistic Regression obtained macro F1 0.7154 and
goodware F1 0.4469, while the BiLSTM obtained macro F1 0.7651 and goodware F1 0.5542. This suggests
that modeling the API order can help the minority class. However, the improvement is moderate given
the much higher model complexity and training cost. The frequency baseline remains very competitive,
which suggests that API-call frequencies already capture a large part of the malware/goodware signal
in this dataset.

\begin{table}[h!]
\centering
\begin{tabular}{lcccc}
\hline
\textbf{Model} & \textbf{Accuracy} & \textbf{Balanced acc.} & \textbf{Macro F1} & \textbf{Goodware F1} \\
\hline
Logistic Regression baseline & 0.9688 & 0.6654 & 0.7154 & 0.4469 \\
Best RNN: BiLSTM             & 0.9545 & 0.8597 & 0.7651 & 0.5542 \\
\hline
\end{tabular}
\caption{Comparison between the frequency-based baseline and the best RNN model on the test set.}
\label{tab:task3_baseline_comparison}
\end{table}
